\documentclass[10pt]{article}

% graphicx package, useful for including eps, pdf, and raster graphics
\usepackage{graphicx}
\DeclareGraphicsExtensions{.png,.png,.jpg,.pdf}

% basic packages
\usepackage{fancyhdr}
\usepackage{color}
\usepackage{parskip}
\usepackage{float}
\usepackage{microtype}
\usepackage{url}
\usepackage{hyperref}
\usepackage{booktabs}
\usepackage{makecell}
\usepackage{multirow}
\usepackage{pbox}
\usepackage{enumitem}

% fonts
\usepackage{fontspec}
\setmainfont{DroidSans}[
  Extension   = .ttf,
  UprightFont = *,
  BoldFont    = *-Bold ]
\setmonofont{DroidSansMono}[
  Extension   = .ttf,
  UprightFont = * ]

\hypersetup{colorlinks,urlcolor=blue}
\urlstyle{same}

% page layout
\usepackage[top=1in, bottom=1in, left=0.75in, right=0.75in, includefoot, heightrounded]{geometry}

% Useful for seeing the layout as you tweak it.
%\usepackage{showframe}

\setlength{\headheight}{1.1in}
\pagestyle{fancy}
\fancyhf{}

\lhead{\includegraphics[height=1in,clip]{logo}}
\rhead{\pbox[b][1in][c]{\textwidth}{
    \small
    \textbf{Baaritaanka waxaa samaynaaya:}\\
    Northwest Genomics Center\\
    University of Washington School of Medicine\\
    Genome Sciences Box 355065\\
    3720 15th Ave NE\\
    Seattle, WA 98195-5065}}

\newcommand{\PageLine}{\rule{\textwidth}{0.25mm}}

\begin{document}

\begin{description}[font=\normalfont,align=left,labelwidth=12em]
\item [\textbf{\VAR{pat_name|e}}]
\item [\textbf{Taariikhda Dhalashada}] \VAR{birth_date|e}
\item [\textbf{Unugta}] Dheecaan laga qaaday bartamaha sanka, oo qofku iska qaadaayo
\item [\textbf{Koodhka saambalka}] \VAR{qrcode|e}
\item [\textbf{Taariikhda aruurinta}] \VAR{collect_ts|e}
\item [\textbf{Taariikhda Warbixinta}]
  %- if status_code == "never-tested"
  Waligay layma baarin
  %- else
  \VAR{result_ts|e}
  %- endif
\end{description}

\PageLine

\begin{center}
\Large
\textbf{NATIIJADA BAARITAANKA – Waxaa lagugu boorinayaa inaad natiijadaan la
  wadaagto dhakhtarkaaga}
\end{center}

\bigskip

\large \underline{\textbf{Xogta Ku saabsan Baaritaankaan}}

\begin{itemize}
\item

  COVID-19 waa \textbf{cudur} uu keeno \textbf{fayras} loo yaqaano SARS-CoV-2.
  Baaritaankaan cilmi baarista ah wuxuu raadinayaa fayraska ku jira saambalka
  laga qaaday sankaaga si loo go'aansho inaad qabto iyo inkale COVID-19. Erayga
  farsamo ee baaritaanka noocaan ah waa ``falcelinta isbadalka hide sidayaasha
  polymerase,'' ama ``RT-PCR'' marka lasoo gaabsho.

\end{itemize}

\bigskip

\large \underline{\textbf{Xasuusino Muhiim ah}}

Kahor intaadan fiirin natiijada baaritaankaaga cilmi baarista, waa inaad akhridaa:

\begin{itemize}
\item

  \textbf{Nooca sambalka} lagaa qaaday iyo \textbf{baaritaanka shaybaadhka} ayaa
  qayb ka ah \textbf{darasada cilmi baarista ah} warbixintaan ayaana qeexaysa
  \textbf{natiijada baaristaada cilmi baarista.}

\item

  Adigga ayaa iska qaaday soona calaamadeeyay (ama qof kale ayaa calaamadeeyay)
  waxaadna kusoo celisay shaybaadhka adoo usoo mariyay adeegga agabka geeya. Ma
  uusan qaadin, mana uusan calaamadayn, ama gaari kuma uusan soo qaadiin
  shaqaale caafimaad oo xirfad leh sida kalkaaliso ama dhakhtar. Arintaan ayaa
  samaayn karta natiijada baaritaankaaga, dhakhtar caafimaad ayaana soo jeedin
  kara in mar kale lagu baaro.

\item

  Inaad ka qaybgasho \textbf{daraasadaan cimi baarista ah} ma buuxinayso booska
  daryeelkaaga caafimaad ee caadiga ah. Haddii aad xanuun badan qabto ama aad
  walaac ka qabto natiijada baaritaankaaga, waa inaad la xariirtaa dhakhtar ama
  daryeele caafimaad.

\end{itemize}

\bigskip
\PageLine

\large \underline{\textbf{Natiijada Baaritaankaaga SARS-CoV-2 (COVID-19) ee
  Cilmi Baarista }}

%- if status_code == "negative"
Lagaama helin (Negative (Ma qabtid))\\
%- elif status_code == "positive"
Waa lagaa helay (Positive (Waad qabtaa))\\
%- elif status_code == "inconclusive"
Jawaabtu ma cadda\\
%- elif status_code == "never-tested"
Lama baari karo\\
%- else
\VAR{0/0} %# poor man's raise
%- endif

\PageLine
\bigskip

\large \underline{\textbf{Fahanka Natiijadaada}}

Shaxdaan waxay ku siinaysaa dhammaan natiijooyinka uu baaritaankaagu muujinkaro
wuxuu qeexayaa waxa mid kasta oo natiijooyinkaan kamid ah kadhigan yahay. Waxaad
daabici kartaa ama iimeel ku diri kartaa \textbf{natiijadaan baaritaanka cilmi}
baarista ah si aad ula wadaagto dhakhtarkaaga ama daryeelahaaga caafimaadka.

Nooc kasta oo natiijada baaritaankaagu yahayba, waa inaad guriga joogtaa aadna
ka fogaataa dadka kale haddii aad jiran tahay. \textbf{Haddii astaamaha
xanuunkaagu sii socdaan, aad xanuun badan dareemayso ama aad ka cabsi qabto
xanuunkaaga, ama aad qabto su'aalo la xariira natiijadaada baaritaanka cilmi
baarista ah, waa inaad la xariirtaa dhakhtar ama daryeele caafimaad.}

\begin{itemize}


\item

  \textbf{If your test says ``Not Detected (Negative)''}, it means our lab did not find the
  virus that causes COVID-19 in the sample that you sent to us. No test is
  perfect, and there is a small chance that you could have the virus even though
  your test is negative.

\item

  \textbf{If your test says ``Detected (Positive)''}, it means our
  lab found the virus that causes COVID-19 in the sample that you sent us. You
  should stay home and avoid being around other people, especially if you are
  feeling sick.

  Your result will be shared with the Washington State and local health
  departments. You do \textbf{not} need to report these results yourself. The
  health department may call you to ask questions about your illness and people
  you may have been around while you were sick.

\item

  \textbf{If your test says ``Inconclusive''}, it means our lab is not certain
  of your test result. No test is perfect, and an inconclusive result happens
  when we find a small amount of SARS-CoV-2 in your sample, but not enough for
  us to say that you have COVID-19.

\item

  \textbf{If your test says ``Unable to be tested''}, it means there was either
  not enough material on your sample for us to test, or there was an error when
  we ran the test. Since we are not able to test your sample, we could not find
  whether or not you have the virus that causes COVID-19.

\end{itemize}

\bigskip

\large \underline{\textbf{Agabyada}}

Haddii aad doonayso inaad xog dheeraad ah ka hesho COVID-19, fadlan akhri xogta
kasoo baxda iskaashatada caafimaadka dadwaynaha:

\begin{itemize}
\item
  Public Health \textemdash Seattle \& King County (Caafimaadka Bulshada
  \textemdash Seattle iyo King County):\\
  \url{https://www.kingcounty.gov/depts/health/covid-19/languages/somali.aspx}

\item

  Washington State Department of Health (Waaxda Caafimaadka ee Gobalka Washington):\\
  \url{https://www.doh.wa.gov/Emergencies/NovelCoronavirusOutbreak2020COVID19/somali}

\item
  U.S. Centers for Disease Control and Prevention (CDC, Xarumaha Maraykanka ee
  Ka hortaga iyo Xakamaynta Cudurka):\\
  \url{https://wwwn.cdc.gov/pubs/other-languages?Sort=Lang%3A%3Aasc&Language=Somali}

\end{itemize}

Haddii aad doonayso inaad qof kula hadasho taleefanka, fadlan wac mid kamid ah
kuwaan soo socda:

\begin{itemize}

\item

  \textbf{Su'aalo caafimaad oo la xariira COVID-19:} Kala xariir King County
  Call Center (Xariirka Wicitaanka King County) inta u dhaxaysa 8AM iyo 7PM oo
  aad ka wacayso \textbf{206-477-3977}

\item

  \textbf{Su'aalaha aan ahayn kuwa caafimaadka ee la xariira COVID-19, ayna ku
  jiraan ku dhaqanka iyo arimaha la xariira ganacsiga:} La xariir King County
  COVID-19 Business and Community Information Line (Laynka Xogta Ganacsiga iyo
  Bulshada ee King County u qaabilsan COVID-19), Isniinta - Jamcada inta u
  dhaxaysa 8:30AM iyo 4:30PM ood ka wacayso \textbf{206-296-1608}

\item

  \textbf{Su'aalo guud oo ku saabsan COVID-19 ee Gobalka Washington:} La xariir
  Washington State Novel Coronavirus Call Center (Laynka Xarunta cusub ee
  Gobalka Washington u qaabilsan Korona Fayraska Cusub) inta u dhaxaysa 6AM iyo
  10PM ood ka wacayso \textbf{1-800-525-0127}

\end{itemize}

\bigskip

\large \underline{\textbf{Qaababka iyo Xadiyada}}

\textbf{Qaababka}\\
Nidaamkaan waxaa loo sameeyay in lagu baaro SARs-CoV-2 ayadoo la adeegsanaayo
nuqulada hide sidaha ee PCR. Nidaamkaan wuxuu adeegsadaa sharqanta qaybaha
fayraska ee ORf1B iyo S wuxuuna sidoo kale raadshaa xakamaynta sharqanta RnaseP
ka dadka. Baaritaankaan waxaa abuurtay oo sifooyinkiisa ay go'aamiso University
of Washington’s Northwest Genomics Center.Shaybaadhka waxaa ruqsada siisay
Xarunta Baaritaanka Caafimaadka u qaabilsan WA Department of Health ayadoo u
aqoonsatay inuu u qalmo fulinta baaritaan caafimaad oo shaybaadh ah kaasoo aad u
balaaran.

\textbf{Xadiyada}

\begin{itemize}

\item

  Natiijooyinkaan waxaa lagu bixiyaa ayagoo qayb ka ah daraasada cilmi baarista.

\item

  Adigga ayaa iska qaaday soona calaamadeeyay (ama qof kale ayaa calaamadeeyay)
  waxaadna kusoo celisay shaybaadhka adoo usoo mariyay adeegga agabka geeya. Ma
  uusan qaadin, mana uusan calaamadayn, ama gaari kuma uusan soo qaadiin
  shaqaale caafimaad oo xirfad leh sida kalkaaliso ama dhakhtar. Arintaan ayaa
  samaayn karta natiijada baaritaankaaga, dhakhtar caafimaad ayaana soo jeedin
  kara in mar kale lagu baaro.

\item

  Sabab la xariirta in warbixintaan lagu saleeyay xog laga helay daraasada cilmi
  baarista ah, natiijooyinkaaga ayaa laga yaabaa inaysan jawaab rasmi ah ka
  bixin COVID-19, oo ah cudurka uu keeno caabuqa SARS-CoV-2.

\end{itemize}

\end{document}
